\section{An Exact Three-Dimensional Finite-Cavity Check}
\label{app:exactcavity}

This check separates the deterministic finite-benefit argument from Monte
Carlo reference error. It is an analytic verification case, not an additional
application benchmark or a validation of the floating-point random-walk law.
Let $D=B_1(0)\subset\R^3$, let $f=1$, and impose zero Dirichlet data on all
boundaries. Remove the concentric ball $\overline B_a(0)$, $0<a<1$, and use
compliance $J=\int u$. The source is restricted to the retained domain; this
check therefore does not belong to the fixed-total-heat design comparison.

\begin{proposition}[Exact benefit and certificate gap]
\label{prop:exactcavity}
Writing $s=|x|$, the original and perforated states are
\begin{align}
 u(s)&=\frac{1-s^2}{6},\\
 v_a(s)&=\frac{1-s^2}{6}-c_a\left(\frac1s-1\right),
 \qquad c_a=\frac{a(1+a)}6,\quad a\le s\le1.
\end{align}
The exact finite benefit $B_a=J(D)-J(D\setminus\overline B_a)$ is
\begin{equation}
 \frac{B_a}{4\pi}
 =\frac{a(1+a)^2(1-a)}{36}+\frac{a^5}{45}.
 \label{eq:exactcavitybenefit}
\end{equation}
The average-trace lower bound is
$L_a=4\pi a(1-a^2)^2/36$, and its gap is
\begin{equation}
 \frac{B_a-L_a}{4\pi}
 =\frac{a^2}{180}\bigl[5(1-a^2)+a(5-a^2)\bigr]>0.
 \label{eq:exactcavitygap}
\end{equation}
Moreover, for the relative capacity $\rho_a=4\pi a/(1-a)$,
\[\lim_{a\downarrow0}(-B_a/\rho_a)=-1/36=-u(0)^2.\]
\end{proposition}

\begin{proof}
\textbf{Step 1: solve the radial equations.}
In three dimensions, $\Delta s^2=6$ and $\Delta(1/s)=0$ for $s>0$.
Thus both displayed states satisfy $-\Delta u=-\Delta v_a=1$ on their
respective domains. The original state is regular at zero and vanishes at
$s=1$. The perforated state vanishes at $s=1$, and
\[
 v_a(a)=\frac{1-a^2}{6}
 -\frac{a(1+a)}6\frac{1-a}{a}=0.
\]
Uniqueness of the Dirichlet problem identifies them as the states.

\textbf{Step 2: integrate the objective directly.}
Polar integration gives $J(D)/(4\pi)=1/45$. On the shell, an antiderivative
of $v_a(s)s^2$ is
\[
 F_a(s)=\left(\frac16+c_a\right)\frac{s^3}{3}
       -\frac{s^5}{30}-\frac{c_as^2}{2}.
\]
Consequently $B_a/(4\pi)=1/45-F_a(1)+F_a(a)$.
Substitution of $c_a=a(1+a)/6$ and collection of powers of $a$ gives
\[
 \frac{B_a}{4\pi}
 =\frac{a+a^2-a^3-a^4}{36}+\frac{a^5}{45},
\]
which is \cref{eq:exactcavitybenefit}.

\textbf{Step 3: verify the difference-energy identity independently.}
Extend $v_a$ by zero inside the cavity and put $w=u-v_a$.
Then $w'(s)=-s/3$ for $s<a$, and $w'(s)=-c_a/s^2$ for $a<s<1$.
It follows that
\begin{align}
 \frac1{4\pi}\int_D|\nabla w|^2
 &=\int_0^a\frac{s^4}{9}\,ds
   +c_a^2\int_a^1s^{-2}\,ds\\
 &=\frac{a^5}{45}+c_a^2(1/a-1)
 =\frac{B_a}{4\pi}.
\end{align}
The contribution inside the removed ball is necessary; dropping it would
not give the exact finite benefit.

\textbf{Step 4: evaluate and compare the trace certificate.}
The trace is constant on the sphere: $\overline u_a=(1-a^2)/6$.
The center Green integral is exactly $a^2/6$, so
$u(0)-a^2/6=\overline u_a$. Thus the exact-state version of
\cref{thm:averagecertificate} returns $L_a=4\pi a\overline u_a^2$.
Subtracting it from Step 2 yields \cref{eq:exactcavitygap}.
Both terms in its square brackets are positive for $0<a<1$.

\textbf{Step 5: check the infinitesimal normalization.}
The relative equilibrium potential equals one for $s\le a$ and
$q_a(s)=a(1/s-1)/(1-a)$ for $a<s<1$.
Integrating its squared radial derivative gives
\[
 \rho_a=4\pi\frac{a^2}{(1-a)^2}\int_a^1s^{-2}\,ds
       =\frac{4\pi a}{1-a}.
\]
Dividing \cref{eq:exactcavitybenefit} by this expression and taking
$a\downarrow0$ gives the claimed $-1/36$ derivative. This also distinguishes
relative capacity from the free-space ball capacity $4\pi a$ at finite radius.
\end{proof}

We check the boundary trace, the two independent expressions for benefit, and
the certificate gap using exact rational arithmetic at
$a=i/1000$, $i=1,\ldots,999$. All algebraic residuals are exactly zero.
The finite-benefit lower bound retains $99.00\%$, $95.00\%$, $89.99\%$,
and $74.80\%$ of the exact benefit at radii $0.01$, $0.05$, $0.10$, and
$0.25$, respectively. These values describe this symmetric model only.
They neither establish certificate tightness on general meshes nor verify
sampling coverage for the implemented validator.
